%%
% Springer LNCS template (match 978-981-97-8798-2 sample layout).
% Compile: pdflatex "main copy 2.tex"; bibtex "main copy 2"; pdflatex x2.
% Requires: llncs.cls, splncs04.bst (TeX Live: texlive-publishers; or Springer/CTAN).
\documentclass[12pt]{llncs}

\usepackage[UTF8]{ctex}
\usepackage[top=2.5cm,bottom=2.5cm,left=2.2cm,right=2.2cm]{geometry}
\usepackage{graphicx}
\graphicspath{{figures/}}
\usepackage{float}
\usepackage{amsmath,amssymb}
\usepackage[square,numbers]{natbib}
\usepackage{enumitem}
\setlist[enumerate]{leftmargin=1.2em, itemindent=0em, labelwidth=1em, labelsep=0.4em}
\setlist[itemize]{leftmargin=1.2em, labelsep=0.4em}
\usepackage[table]{xcolor}
\usepackage{booktabs}
\usepackage{array}
\usepackage{multirow}
\usepackage{hyperref}
\hypersetup{colorlinks=true, linkcolor=blue, citecolor=blue, urlcolor=blue}
\usepackage{caption}
\captionsetup[figure]{aboveskip=4pt, belowskip=4pt}
\captionsetup[table]{aboveskip=4pt, belowskip=4pt}
\setlength{\textfloatsep}{6pt plus 2pt minus 1pt}
\setlength{\intextsep}{6pt plus 2pt minus 1pt}

\begin{document}
\title{CIPHERNOVA：一种基于测试反馈与分层提示的密码学代码生成方法}
\titlerunning{CIPHERNOVA：测试反馈与分层提示驱动的密码学代码生成}

% ===== Anonymous switch =====
\newif\ifanonymized
\anonymizedtrue      % 投稿匿名版
% \anonymizedfalse   % 终稿实名版

\ifanonymized
\author{Anonymous Authors}
\institute{Anonymous Institution\\
\email{anonymous@anonymous.edu}}
\authorrunning{Anonymous Authors}
\else
\author{Jia-Yi He\inst{1} \and Xiao Lan\thanks{Corresponding author: lanxiao@scu.edu.cn}}
\institute{School of Cyber Science and Engineering, Sichuan University, Chengdu, Sichuan, China\\
\email{2022141500268@stu.scu.edu.cn, lanxiao@scu.edu.cn}}
\authorrunning{J.\ He and X.\ Lan}
\fi
\maketitle

\begin{abstract}
传统大语言模型（LLM）在密码学代码生成中普遍面临错误率高、缺乏领域优化等问题。本文提出一种系统化方法，包含五项核心设计：（1）基于测试反馈的自适应改进机制——闭环验证、标准测试向量、自动错误模式识别与针对性修复指导；（2）分层可组合的领域提示工程方法——四层架构与 101 个精炼模板，并采用\emph{失败驱动的提示精炼}；（3）密码学错误分类框架，含四大类及专用检测与修复策略；（4）兼容离线本地模型（如 Qwen2.5-Coder:7B）；（5）\textbf{面向本地模型的教师示例检索注入}——以通过标准测试的云端 LLM 为教师，在生成时注入少样本通过示例、在改进时注入同题参考代码，在不重训练前提下提升本地模型功能测试通过率。在 159 个真实测试用例的受控评测设置下，云端 LLM 在本基准上达到 100\% 的\emph{代码生成与提取成功率}与 100\% 的\emph{功能测试通过率}；本地 7B 模型也达到 30\% 的\emph{功能测试通过率}，但生成时延显著高于云端模型。本文方法并在同口径子集上明显优于通用单提示基线，并且\textbf{由反复失败驱动的持续提示优化}是取得上述结果最耗时且最关键的一环。本方法为高可信密码学代码的自动化生成提供了系统化解决方案，但实验结论适用范围限于本文基准与评测设置。
\end{abstract}

\keywords{密码学代码生成 \and 大语言模型 \and 测试反馈 \and 提示工程 \and 错误模式分类 \and 代码自改进}

\section{引言}
\subsection{研究背景与问题陈述}
随着网络空间安全威胁日益加剧，密码算法已成为数据加密、身份认证与隐私保护等关键场景的核心支撑技术。然而，密码代码的实现对专业性要求极高：一方面涉及 Feistel 网络、公钥密码数论基础等复杂数学原理；另一方面必须严格遵循 NIST、FIPS 等国际标准。任何细微的实现偏差（如填充模式错误、IV 使用不当）都可能导致整个系统的安全崩溃。

对于非密码学专长的开发者而言，编写高可信的密码代码面临很大挑战。查阅标准文档、复用开源库等传统方式存在三大不足：学习成本高（密码标准文档往往晦涩难懂）、灵活性差（开源库难以适配个性化业务）、安全风险高（开源代码可能含有未发现漏洞）。

近年来，由 GitHub Copilot、CodeT5 等大语言模型（LLM）驱动的代码生成技术为上述问题提供了新思路。然而，通用 LLM 在密码学代码生成任务上存在明显局限：一是错误率居高不下，密码代码对填充方式、初始向量、密钥编排等细节要求严格，通用模型易产生理解偏差导致功能失败；二是提示策略缺乏领域针对性，现有通用提示模板无法覆盖密码学的标准符合性、安全约束与算法细节等特殊需求；三是纠错困难，生成代码测试失败时往往依赖人工干预；四是评估体系不完善，Pass@k 等通用代码生成指标仅关注语法正确性，无法衡量密码代码的安全性与标准符合性。

因此，如何面向密码学领域为 LLM 设计代码生成优化方法，以提升代码正确性、安全性与标准符合性，已成为亟待解决的重要问题。

\subsection{错误密码代码的安全影响}
\label{subsec:security_implications}
错误或不符合标准的密码实现会直接削弱部署时的安全性。现实事件与 CVE 表明，固定或可预测 IV、ECB 模式误用、弱密钥派生或错误填充等（如 mycli 的 CVE-2023-44690、常用库中的 API 误用）可导致机密性或完整性的完全沦陷。因此，降低此类错误代码的生成与部署比例是一项具体的\emph{安全保证}贡献：本方法通过将测试反馈与领域提示、错误模式驱动的修复相结合，我们将标准从“能编译的代码”提升为“通过标准测试向量且满足密码学约束的代码”，从而减少 LLM 生成的密码代码在安全敏感场景中引入的攻击面。

\subsection{研究目标与主要贡献}
针对上述问题，本文提出一种基于测试反馈与分层提示的密码学代码生成方法。核心目标是通过测试反馈的闭环改进与分层提示的精准引导，提升 LLM 生成的密码学代码质量，减轻领域开发者的负担。

本文的主要贡献可概括为：
\begin{itemize}
    \item 提出\textbf{基于密码学代码错误模式分类的测试反馈与自适应改进机制}。建立四类错误的分类框架并实现六步闭环流程；通过分析测试输出自动识别错误类型并生成针对性修复指导。
    \item 设计\textbf{组合领域提示工程方法}。构建四层提示架构（通用 → 算法 → LLM → 语言模板），形成 101 个精炼模板，并支持\textbf{失败驱动的提示精炼}——通过迭代强化提示以修复反复出现的生成失败（如缺失 \texttt{stdint.h}、未定义引用、将明文当密文输出、输出格式错误等）。该精炼过程被证明是实验工作中最具挑战且最耗时的部分，对报告的通过率至关重要。
    \item 研制兼容\textbf{离线本地大模型}的密码学代码生成引擎。兼容 Qwen2.5-Coder:7B 等轻量级本地大模型，在不依赖公网的前提下保证生成效率、正确性与数据隐私安全。
    \item 引入\textbf{面向本地模型的教师示例检索注入}：以通过标准测试的云端 LLM 为教师，在生成时注入同（算法、模式、语言）的少样本通过示例，在改进时注入同题参考代码，在不重训练的前提下提升本地模型在支持算法（如 SM4）上的功能测试通过率。
\end{itemize}
实验构建了包含 159 个真实测试用例的密码学代码生成数据集，覆盖 4 种主流算法与 3 种编程语言，为相关研究提供了标准化评估基准。通过对比实验与消融实验，验证了所提方法的有效性及各模块的贡献。


\subsection{论文组织结构}
本文后续章节安排如下：第 2 节综述相关研究；第 3 节阐述方法的整体架构及三个核心创新模块的设计细节；第 4 节通过实验验证方法性能；第 5 节总结全文并讨论局限与未来方向。

\section{相关工作}
本节从大语言模型代码生成、领域代码生成、提示工程与代码自改进技术四个维度综述相关研究，并在各小节末简要说明本文工作如何受其启发或在其基础上拓展。

\subsection{大语言模型代码生成}
大语言模型（LLM）在代码生成上的发展依赖大规模代码预训练与自回归生成。CodeGen~\cite{codegen1}、CodeGen2~\cite{codegen2} 建立了开放代码语言模型及规模与质量的经验；Chen 等~\cite{codex} 对代码 LLM 与代码生成能力进行了评估。CodeT5~\cite{codet5}、CodeT5+~\cite{codet5plus} 提供编码器-解码器代码模型；InCoder~\cite{incoder} 侧重代码填空与综合；StarCoder~\cite{starcoder}、WizardCoder~\cite{wizardcoder}、Code Llama~\cite{codellama}、DeepSeek-Coder~\cite{deepseekcoder} 等在 HumanEval 等基准上取得突出结果。CodeBERT~\cite{codebert}、GraphCodeBERT~\cite{graphcodebert} 面向代码表示进行预训练；CodeXGLUE~\cite{codexglue} 提供代码理解与生成基准。上述工作主要面向通用代码任务，语料中密码学内容少且缺乏正确性与标准符合性的专门约束，难以直接用于密码学代码。\textbf{本文以这类 LLM 为骨干}（云端 API 与本地模型），\textbf{通过领域提示与测试反馈闭环进行增强}，使其在不重训练的前提下可靠生成符合标准的密码学代码。

\subsection{领域代码生成}
将领域知识融入生成过程是提升专业代码质量的关键。在密码学方向，Li 等~\cite{encgpt} 提出 EncGPT，采用多智能体工作流进行动态加密算法生成；Baek 等~\cite{cryptollm} 在 ESORICS 2024 上探索了基于提示的 LLM 密码 API 误用检测。两者在密码学代码生成上做了有益尝试，但均缺乏系统的错误模式分析与基于测试反馈的闭环改进机制。Zhang 等~\cite{seccoder} 在 EMNLP 2024 提出 SecCoder，面向可泛化、鲁棒的安全代码生成，验证了领域知识注入对安全代码生成的有效性。总体而言，在密码正确性、NIST/OpenSSL 等标准符合性、错误模式分类以及“生成-验证-改进”闭环方面仍存在研究空白。\textbf{受 EncGPT 的领域结构运用与 SecCoder 中领域注入提升正确性的启发，我们设计了密码学错误模式分类并将其类别注入生成提示与修复指导}，使模型获得显式的、按算法与语言区分的约束而非通用指令。

\subsection{提示工程}
高质量提示能提升 LLM 在推理与代码任务上的表现。Chain-of-Thought (CoT)~\cite{cot}、零样本 CoT~\cite{zeroshotcot}、Plan-and-Solve~\cite{planandsolve} 改进推理；Brown 等~\cite{fewshot}、Austin 等~\cite{progsynprompt} 表明少样本与结构化提示有助于程序综合。Wang 等~\cite{coderagbench} 探索了代码生成的检索；Shrivastava 等~\cite{repolevelprompt} 研究了仓库级提示生成。Budagam 等~\cite{hierprompt} 提出分层提示分类，Madaan 等~\cite{selfrefine} 利用迭代自反馈（Self-Refine）精炼输出。上述工作均未针对密码学设计可随算法、语言与模型差异适配的分层可组合提示架构。\textbf{我们借鉴 Budagam 等~\cite{hierprompt} 的分层可组合提示与 Self-Refine~\cite{selfrefine} 的迭代精炼，在密码学场景下实例化为}四层（通用、算法、LLM、语言）及失败驱动的精炼，使每种（算法、语言、模型）组合获得定制化约束。

\subsection{代码自改进技术}
“生成-验证-改进”闭环是提升代码质量的有效范式。Self-Edit~\cite{selfedit} 提出面向生成的故障感知代码编辑器；LeDeX~\cite{ledex} 训练 LLM 自调试与解释代码；Self-Refine~\cite{selfrefine} 利用自反馈进行迭代精炼；RefineCoder~\cite{refinecoder} 通过自适应批评精炼改进代码生成；AgentCoder~\cite{agentcoder} 采用多智能体迭代测试与优化。CYCLE~\cite{cycle}、RLEF~\cite{rlef} 利用执行反馈；RepairAgent~\cite{repairagent} 提出基于 LLM 的程序修复智能体；Li 等~\cite{li2023} 将 LLM 与程序分析结合用于自动修复；RGD~\cite{rgd}、UniDebugger~\cite{unidebugger} 采用多智能体或分层调试；Tang 等~\cite{tang2024} 研究了 LLM 代码修复中的探索-利用。ReCode~\cite{recode} 用检索改进 LLM 代码修复。在安全方向，Basic 等~\cite{secsurvey} 综述了 LLM 在代码安全中的应用，He 等~\cite{securecodegen} 研究了安全代码生成的指令微调。上述工作多聚焦通用代码、测试生成或安全代码，缺乏面向密码学代码的、由测试反馈驱动的错误模式分类与针对性修复策略。\textbf{遵循 Self-Edit~\cite{selfedit}、RefineCoder~\cite{refinecoder}、AgentCoder~\cite{agentcoder} 的“生成-验证-改进”范式，我们实现了六步闭环}（生成→标准测试向量验证→错误分析→修复指导→改进→迭代），但修复步骤由\emph{密码学}错误模式分类器与领域修复模板驱动，而非通用修复提示，从而将 IV 误用、输出格式错误等失败映射为具体的提示约束。

跨论文对比前须对齐协议（测试桩、迭代轮数、通过准则）；通用代码基准侧重单次合成与 Pass@k，与本文标准向量闭环设定不同。

\subsection{研究空白小结}
\label{subsec:research_gaps}

据我们所知，尚无工作将\emph{测试反馈驱动的闭环}、面向密码学的\emph{分层可组合提示架构}与\emph{系统化的密码学错误模式分类}统一于一个密码学代码生成框架。现有研究存在三方面空白：（1）缺乏面向密码学代码的系统化错误模式分类框架；（2）未设计可随算法、语言与 LLM 适配的分层可组合提示方法；（3）未构建基于测试反馈的密码学代码自改进闭环。本文针对上述空白提出系统化解决方案。

\section{方法设计}
基于上述研究空白，并依托通用代码 LLM 作为骨干、领域注入以保障正确性（如 EncGPT~\cite{encgpt}、SecCoder~\cite{seccoder}）、分层与迭代提示设计~\cite{hierprompt,selfrefine} 以及“生成-验证-改进”闭环~\cite{selfedit,refinecoder,agentcoder}，本节给出本文方法。阐述基于测试反馈与分层提示的密码学代码生成思路，整体框架见图~\ref{fig:framework}。系统采用五层架构：\emph{用户交互层}（Web UI 与 CLI）；\emph{API 服务层}（FastAPI，提供生成、改进与连接测试等接口）； \emph{核心业务层}，\emph{CryptoAgent} 负责分层提示拼装、调用 LLM、代码提取与清洗及验证桩执行，\emph{LLMAdapter} 对接云端 API 与本地 Qwen2.5-Coder:7B，图中 \textbf{核心模块 A}（基于测试反馈的自适应改进）落在该层并在验证失败时触发改进闭环；自上而下第四层为 \emph{工具支撑层}，包含配置管理、验证引擎、提示加载器、模型评估器，并在该层实现 \textbf{核心模块 B}（分层可组合领域提示）与 \textbf{核心模块 C}（密码学错误模式分类，为上层错误分析提供类别信号）；第五层（最下）为 \emph{外部服务层}（云端LLM与本地模型）。

\begin{figure}[htbp] 
  \centering
  \includegraphics[width=\textwidth]{Framework.png} 
  \caption{基于测试反馈与分层提示的密码学代码生成系统整体框架。五层自上而下为用户交互、API 服务、核心业务、工具支撑、外部服务。}
  \label{fig:framework}
\end{figure}
本方法主要由三个核心模块构成：基于测试反馈的自适应改进机制、分层可组合的领域提示工程方法，以及密码学代码错误模式分类系统。


\subsection{基于测试反馈的自适应改进机制}
基于测试反馈的自适应改进机制是本文方法的核心，旨在建立“生成、验证、分析、修复、改进、迭代”的闭环改进流程，实现代码质量的持续提升。该机制包含三个关键步骤：错误模式识别、测试反馈分析与针对性修复指导生成。

\subsubsection{问题分析。}
传统代码生成系统采用线性流程，测试失败时仅简单重试，不分析错误根因也不做针对性优化。密码学代码错误涉及算法逻辑、实现细节、环境配置等多维因素，比通用代码复杂得多，单纯重试无法解决根本问题；现有通用修复策略也缺乏密码学专业支撑，难以适应填充模式调整、初始向量管理等特殊修复需求。

\subsubsection{方法设计。}
自适应改进机制包含三个核心组件：错误模式识别、测试反馈分析与针对性修复。各组件设计要点如下：
\begin{enumerate}
    \item \textbf{错误模式识别} \\
    目标是根据测试输出准确定位代码中的错误类型。结合密码学代码特点，本文设计了多级错误检测策略与基于规则的检测器，通过针对性函数封装实现错误特征的快速识别。检测从输出完整性、内容一致性、算法完整性与占位符存在性四个维度进行，并辅以语法错误检测与安全约束检测，形成综合检测体系。

    \item \textbf{测试反馈分析} \\
    目标是从测试输出中提取关键错误信息，为生成修复指导提供精确依据。本文设计的测试反馈分析模块通过五步流程提取关键信息，包括输出缺失检测、长度差异分析、明文输出检测、导入缺失检测与错误模式汇总，为修复指导提供精确依据。

    \item \textbf{针对性修复指导生成} \\
    目标是根据错误特征向量生成领域化、可操作性强的修复指令。修复指导采用结构化指令格式，统一包含错误类型、修复要求与示例代码三部分，根据识别出的错误模式生成精确指导方案，既明确修复方向又提供可参考示例，降低 LLM 修复时的理解偏差。修复指导生成模块采用“模板匹配 + 规则生成”的融合方式，为不同错误类型预设基础模板，并根据错误特征向量中的具体信息（如算法类型、编程语言、错误细节）动态补充内容，保证指导的精确性与可操作性。
\end{enumerate}

\subsubsection{闭环改进流程。}
基于上述三个核心组件，本文构建了“生成、验证、分析、修复、改进、迭代”的六步闭环改进机制，具体步骤如图 \ref{fig:loop-process} 所示：
\begin{enumerate}
    \item \textbf{初始生成}。向 LLM 输入分层组合提示，生成初始代码；
    \item \textbf{代码验证}。使用 NIST、FIPS 等标准测试向量进行自动化测试，收集实际输出、错误信息等反馈数据；
    \item \textbf{错误分析}。通过错误模式识别确定错误类型，基于测试反馈分析生成错误特征向量，定位核心错误点；
    \item \textbf{修复指导生成}。根据错误特征向量输出结构化的针对性修复指令；
    \item \textbf{代码改进}。结合修复指令引导 LLM 生成改进代码；
    \item \textbf{迭代终止}。若改进代码通过测试或达到最大迭代次数（3 次）则终止；否则回到第 2 步继续迭代。
\end{enumerate}
\begin{figure}[htbp]
  \centering
  \includegraphics[width=0.85\textwidth]{SixStepIteration2.pdf}
  \caption{六步闭环改进流程}
  \label{fig:loop-process}
\end{figure}
以 DES-CFB 等模式为例，闭环依次完成初始生成、标准向量验证、错误分析、修复指导注入与代码改进；未通过则重复后三步直至上限 3 轮。
\subsection{分层可组合的领域提示工程方法}
分层可组合的领域提示工程方法旨在通过设计分层提示模板，对 LLM 实现精准引导。本文提出的提示架构由四层组成：通用模板层、算法模板层、LLM 模板层与语言模板层，每层包含多个精炼子模板。

\subsubsection{分层提示架构设计。}
针对代码生成质量不稳定、验证通过率低、跨场景与跨模型兼容性差等问题，我们提出分层提示架构设计。该多层框架将复杂的提示工程需求分解到不同粒度，便于对通用规范、算法特性、模型特点与语言特性进行针对性优化。四层核心结构如下：

\begin{enumerate}
    \item \textbf{通用模板层}。定义与场景无关的内容（系统角色、代码规范、输入/输出格式及基线要求）。统一输出格式显著提高了代码提取准确率。

    \item \textbf{LLM 模板层}。针对不同模型特性（大规模与本地等）进行优化，按提供商与 \texttt{llm.yaml} 等侧载模板差异化拼装，减少跨场景错误。

    \item \textbf{算法模板层}。聚焦算法特有属性（原理、核心步骤、参数要求）。引入按算法定制的提示有效提升了代码验证通过率。

    \item \textbf{语言模板层}。按编程语言明确密码库调用规则、代码风格与防错指南，该层明显增强了代码可编译性。
\end{enumerate}

\subsubsection{失败驱动的提示精炼。}
在上述四层架构基础上，我们强调实践中不可或缺的\emph{由反复生成失败驱动的迭代提示精炼}。

\textbf{失败驱动的提示精炼。}要达到高功能通过率，需在算法–语言–LLM 组合上大量迭代。实验中高频失败包括：缺 \texttt{stdint.h}/辅助函数未定义导致的编译链接错误；明文当密文输出、\texttt{\%04x} 与字节输出不符；IV 来源或输出多余 IV；错误密码库或 API；以及 \texttt{goto cleanup}、C++ 字符串判断等结构问题。
针对每种（算法、语言、LLM）组合，我们迭代精炼提示并加入显式约束：强制头文件与函数定义顺序（辅助函数在 \texttt{main()} 之前）、按字节输出密文（\texttt{printf("\%02x", ...)}）、CBC 正确性（明文与 IV 异或后再块加密，禁止输出 IV 或明文）、缓冲区大小、Python \texttt{import os} 与正确库（如 SM4 用 gmssl、DES 用 pycryptodome）、IV 仅来自 \texttt{TEST\_IV}，以及 C++ 字符串/控制流规则（如用 \texttt{.empty()} 而非 \texttt{!}）。我们还引入轻量后处理（如验证阶段在需要时自动插入 \texttt{\#include <stdint.h>}；确保 C 代码在编译前经过缺头文件修复）。该失败驱动的提示精炼是流程中\textbf{最耗时且最具挑战}的部分，对报告的 FTPR 至关重要。

\subsubsection{各阶段的具体挑战与提示对策。}
\label{subsec:challenges_prompt_table}
下文系统说明生成与改进两阶段的挑战与提示级对策。\textbf{生成}阶段依托分层：算法与模型相关约束（强制头文件、单文件自包含、按字节密文输出、CBC 语义、缓冲区边界、输入接受）以及针对脆弱（算法、语言、模型）三元组的专用错误/正确示例与显式禁止（如密文禁止 \texttt{\%08x}/\texttt{\%04x}）。输出不完整或截断时触发\emph{重试}提示，重新注入上次错误并重申头文件、控制流与截断规则（如 \texttt{\& 0xFF} 不得退化为 \texttt{\& 0}）。\textbf{改进}阶段组合失败代码片段、测试反馈（实际与预期、编译日志）及按优先级排列的失败摘要，提供精确的 \texttt{\#include} 映射、CBC/IV 规范与 API 签名（\texttt{BN\_hex2bn}、\texttt{EVP\_PKEY\_encrypt}）。两阶段共享同一套失败驱动的约束用语。

通用层强制优先使用 OpenSSL、OpenSSL 3.0 下 DES 的 legacy-provider 加载以及严格的不截断规范。模型特定层提供语言规则（C：禁止 \texttt{goto cleanup}，所有清理在 \texttt{return 0} 前完成；C++：\texttt{std::string} 用 \texttt{.empty()} 或指针检查，不用 \texttt{!}；\texttt{vector}、\texttt{cstring} 所需头文件）以及重申头文件、控制流与完整表达式的重试文本。算法-语言层编码模式相关修复：如 SM4-CBC C（固定 include 顺序、单文件辅助函数顺序、仅密文 \texttt{\%02x} 输出）；DES-CBC Python（IV 来自环境变量，禁止默认/零；pycryptodome）；SM4 Python（仅 gmssl）；RSA C/C++（正确使用 \texttt{BN\_hex2bn}/\texttt{EVP\_PKEY\_encrypt}，不用 \texttt{RSA\_public\_encrypt}）；DES-ECB（仅首块输出）；AES-GCM（密文+tag，AAD 来自环境变量）；AES-CFB（CFB-8，非 CFB-128/OFB）。观察到的失败被系统化地转化为 101 个精炼模板中的显式约束。

分层架构与失败驱动精炼相结合，系统性地应对实践中反复出现的错误，提升了跨算法、语言与 LLM 的稳定性与可靠性。

\subsection{面向本地模型的教师示例注入（检索式蒸馏）}
\label{subsec:distillation}
为在不重训练的前提下提高离线本地模型（如 Qwen2.5-Coder:7B）的功能测试通过率，我们以通过标准测试的云端 LLM 产出为\textbf{教师}，将已通过验证的实现检索注入提示：检索键按算法、工作模式、编程语言与操作类型匹配；可将错误簇并入键以实现按失败类型的教师对齐；并对教师样本做质量与白名单门控及 provider 权重排序。改进阶段将测试输出、可选修复摘要、教师参考块与学生待改代码按模板顺序拼装；系统实现为\textbf{检索—注入}，不进行梯度训练。

\textbf{首轮生成。}使用本地模型时，在提示前注入一至多条少样本示例，来源为数据集中同（算法、模式、语言）的通过代码，以减少缺头文件、错误 API 等反复错误。

\textbf{改进阶段。}若本地代码未通过测试，除测试反馈外可注入同条件的教师参考实现，对 SM4-CBC C 等易错题可提高改进轮次成功率。

\textbf{数据收集与开关。}教师数据可由云端模型通过后自动写入 JSONL，或从历史记录回填；路径与最大少样本条数由配置指定（如 \texttt{distillation.dataset\_path}、\texttt{max\_few\_shot\_examples}）。仅当启用该机制且当前为本地提供方时注入；云端独占流程不受影响。

\textbf{工程增强。}支持白名单与最小质量分阈值、按 provider 动态调权，以及改进阶段双教师时在实现栈标签上的差异化选取，以提高对比信息量、避免近重复占上下文。

在 45 个（算法、模式、语言）组合（四算法全模式 × 三语言）上，本地模型启用教师注入前后的整体功能测试通过率见表~\ref{tab:distillation_overall}；蒸馏后 100.0\% 仅为该受控集上的观测值。
\begin{table}[H]
  \centering
  \footnotesize
  \caption{教师示例注入前后整体通过率（本地模型；45 项）。}
  \label{tab:distillation_overall}
  \begin{tabular}{l c c}
    \toprule
    & 通过数 & 通过率 \\
    \midrule
    注入前整体通过率 & 37 / 45 & 82.2\% \\
    注入后整体通过率 & 45 / 45 & 100.0\% \\
    提升             & +8 项   & +17.8\% \\
    \bottomrule
  \end{tabular}
\end{table}

\subsection{密码学代码错误模式分类系统}
\subsubsection{错误模式分类。}
密码学代码错误模式分类系统是本文方法的理论基础，旨在为错误识别与纠正提供系统化框架。通过对 159 个真实测试用例的错误类型分析，本文构建了包含四大类及若干子类的错误分类体系。表 \ref{tab:error_classification} 概括了四类及其代表性示例；更细化的分层约束与失败驱动对策见第~\ref{subsec:challenges_prompt_table}~节及上文。
\begin{table}[htbp]
  \centering
  \footnotesize
  \caption{密码学代码错误模式分类（四类）}
  \label{tab:error_classification}
  \setlength{\tabcolsep}{6pt}
  \begin{tabular}{p{0.3\textwidth}p{0.65\textwidth}}
    \toprule
    大类 & 代表性示例 \\
    \midrule
    算法理解错误 & 模式混淆（如 OFB-64 与 OFB-8）、填充/密钥编排错误。 \\
    实现细节错误 & IV 使用（默认/随机或输出中含 IV）、填充处理、字节序。 \\
    环境配置错误 & 环境变量命名（KEY 与 TEST\_KEY）、缺库、OpenSSL 版本。 \\
    代码结构错误 & \texttt{goto cleanup} 无效、缺头文件、辅助函数不完整（未定义引用）。 \\
    \bottomrule
  \end{tabular}
\end{table}

\subsubsection{修复策略设计。}
针对每类错误，本文设计专用修复策略，保证修复方案的针对性与有效性，核心思路是根据错误类型提供明确的修复要求与示例代码。

算法理解错误修复侧重补充专业知识：密钥长度错误需明确 DES 密钥的标准长度并指导截取前 8 字节；OFB 模式错误需指定使用 OFB64 模式并推荐 EVP\_des\_ofb()；填充模式错误需说明库的自动填充机制并指导删除手动填充代码；密钥编排错误需提供正确置换表的定义与使用示例。

\subsubsection{错误模式匹配算法。}
错误模式匹配采用精确匹配、模式匹配与启发式匹配三种多级策略，协同实现错误模式的快速、精确识别，为生成修复指导提供可靠依据。

\section{实验与结果分析}
本节通过对比实验、消融实验与案例分析验证所提方法的有效性，从多维度系统评估密码学代码生成的性能。

\subsection{实验设置}
\subsubsection{数据集构建。}
为保证实验评估的真实性与代表性，本文构建了包含 159 个真实测试用例的密码学代码生成数据集，其统计特征见表~\ref{tab:dataset_stats}。这些用例覆盖 4 种主流密码算法与 3 种编程语言（C/C++/Python），均使用 NIST、FIPS 及中国国密标准（GM/T）规定的标准测试向量进行验证。

\subsubsection{可复现性。}
为支持可复现性，本文提供代码与 159 个测试用例数据集，并提供包含代码库、提示模板与评估脚本的匿名仓库；数据集与标准测试向量的描述见本文，可根据所引 NIST、FIPS、GM/T 标准复现。

\begin{table}[htbp]
  \centering
  %\footnotesize
  \caption{密码学代码生成数据集统计}
  \label{tab:dataset_stats}
  \setlength{\tabcolsep}{3pt}
  \begin{tabular}{p{0.12\textwidth} 
    >{\centering\arraybackslash}p{0.16\textwidth}  
    >{\centering\arraybackslash}p{0.3\textwidth}  
    >{\centering\arraybackslash}p{0.3\textwidth}  
    }
    \toprule
    算法 & 测试用例数 & 语言分布 (C/C++/Python) & 标准测试向量来源 \\
    \midrule
    DES       & 70                  & 21/24/25                            & FIPS 46-3                      \\
    AES       & 71                  & 26/25/20                            & FIPS 197                       \\
    SM4       & 15                  & 7/4/4                              & GM/T 0002-2012                 \\
    RSA       & 3                   & 1/1/1                              & FIPS 186-4                     \\
    \midrule
    合计     & 159                 & 55/54/50                            & --                             \\
    \bottomrule
  \end{tabular}
\end{table}

\subsubsection{基线系统选择。}
为验证所提方法的有效性与相对优势，本文设置三类对照：\emph{通用外部基线}、\emph{代表性相关工作对比}与\emph{自身模块的消融}。

\textbf{外部基线（通用设置）。}使用相同云端 LLM，仅采用\textbf{单一通用提示}（无算法/语言/LLM 特定模板层，无测试反馈或闭环改进），模拟无任何领域设计的通用代码 LLM 的典型用法。如表~\ref{tab:ablation_prompt} 所示，该基线在\textbf{与下文两张消融表相同}的 DeepSeek、豆包与 Claude ×四算法全支持模式×三语言全网格上仅达 25.0\% FTPR，而四层可组合提示在同一设置下达到 100.0\% FTPR，差距显著，说明分层提示与测试反馈驱动的改进对密码学代码的必要性。

\textbf{相关工作对比。}除通用外部基线外，本文还选取 SecCoder ACL、SVEN LM、Self-Refine 与 AgentCoder 等代表性方法进行对比，并在后文的表~\ref{tab:rw_same_protocol_des} 中给出同一 DES 子集口径下的对齐结果，以比较不同方法在密码学代码生成任务中的实际表现。

\textbf{消融基线}（各自去掉所提方法的一个核心模块以隔离组件贡献）：
\begin{enumerate}[leftmargin=1.2em, labelsep=0.4em]
  \item \textbf{无测试反馈的所提方法}。关闭 Self-Refine 与基于向量测试失败的改进闭环；剥离下一轮生成注入的 \texttt{[VECTOR\_TEST\_RETRY]} 摘要；并关闭 C/C++ 写盘前基于典型失败模式的规则修补链（代码开关 \texttt{\_ablation\_no\_test\_feedback}，覆盖原单独消融的「错误检索自动修复」）。
  \item \textbf{无分层提示的所提方法（\texttt{common\_only}）}。仅保留 \texttt{base\_prompt.yaml} 通用基座，不使用算法/语言等分层模板。
  \item \textbf{无任何提示的所提方法（\texttt{no\_prompt}）}。不加载分层领域模板；用户消息**仅含单行任务描述**（算法、模式、语言与操作），不含通用基座与算法/LLM 层等。
\end{enumerate}

\textbf{消融与错误统计口径。}表~\ref{tab:ablation_main} 与表~\ref{tab:ablation_prompt} 均在\textbf{DeepSeek、豆包（Doubao）与 Claude（Anthropic）}三家云端 LLM 上、于\textbf{四算法全支持工作模式$\times$三语言}的完整任务网格上汇总（VPR/FTPR 为跨三家提供商、该全网格上的聚合结果）。表~\ref{tab:error_repair} 的四类分布与「修复后通过率」则仅基于\textbf{DeepSeek} 在上述\textbf{全算法全模式}网格上写入的 \texttt{llm\_performance.json} 失败记录聚合，\textbf{不含豆包与 Claude}轨迹。

\subsubsection{评估指标定义。}
为全面衡量密码学代码生成质量，定义六项评估指标，涵盖正确性、效率与修复能力：
\begin{itemize}[leftmargin=1.2em, labelsep=0.4em]
  \item \textbf{生成成功率（GSR）}：成功生成可运行代码的测试用例比例；
  \item \textbf{验证通过率（VPR）}：通过语法检查与编译的测试用例比例；
  \item \textbf{功能测试通过率（FTPR）}：输出与标准测试向量一致的测试用例比例；
  \item \textbf{错误修复成功率（ERSR）}：改进后通过功能测试的错误用例比例；
  \item \textbf{平均生成时间（AGT）}：每测试用例的平均生成时间（单位：秒）。
\end{itemize}

\subsection{对比实验结果与分析}
\subsubsection{整体性能对比。}
表 \ref{tab:llm_performance} 给出了所提方法在不同 LLM 上的整体性能对比，主要结论如下：

\begin{table}[H]
  \centering
  \footnotesize
  \caption{所提方法在不同 LLM 上的整体性能}
  \label{tab:llm_performance}
  \setlength{\tabcolsep}{4pt}
  \begin{tabular}{
    p{0.16\textwidth}               
    >{\centering\arraybackslash}p{0.17\textwidth}  
    >{\centering\arraybackslash}p{0.18\textwidth}  
    >{\centering\arraybackslash}p{0.06\textwidth}  
    >{\centering\arraybackslash}p{0.06\textwidth}  
    >{\centering\arraybackslash}p{0.07\textwidth} 
    >{\centering\arraybackslash}p{0.07\textwidth}  
    >{\centering\arraybackslash}p{0.08\textwidth} 
  }
    \toprule
    LLM               & 提示长度 (Python/C/C++) & 生成时间 (Python/C/C++) (s) & GSR (\%) & VPR (\%) & FTPR (\%) & ERSR (\%) & AGT (s) \\
    \midrule
    OpenAI            & 17542/28952/22986            & 37.1/56.6/67.5                     & 100.0    & 100.0    & 100.0     & 81.82     & 178.96  \\
    Deepseek          & 8810/17227/13345             & 21.6/52.0/33.5                     & 100.0    & 100.0    & 100.0     & 59.79     & 56.31   \\
    Doubao            & 13801/24722/19702            & 68.2/158.6/94.3                    & 100.0    & 100.0    & 100.0     & 60.00     & 147.20  \\
    Claude            & 8334/15959/12139             & 31.2/68.0/63.4                     & 100.0    & 100.0    & 100.0     & 30.23     & 34.23   \\
    CodeX             & 17542/28952/22986            & 201.8/336.2/275.5                  & 100.0    & 100.0    & 97.78     & 0.00      & 287.80  \\
    qwen2.5-coder:7b  & 13188/20888/16646            & 2856.0/2857.6/2378.4               & 100.0    & 100.0    & 88.89     & 0.00      & 3720.70 \\
    \bottomrule
  \end{tabular}
\end{table}

\begin{itemize}[leftmargin=1.2em, labelsep=0.4em]
  \item 在本文受控评测设置（159 例、标准向量、固定测试桩、有限迭代预算）下，所提方法在云端 LLM 上达到高且稳定的 GSR/VPR（及多数模型的 FTPR），说明领域提示工程与测试反馈闭环能显著弥补通用 LLM 在密码学知识上的不足；表中数值为该统计口径下的实验结果（FTPR/ERSR 依配置网格与历史合并定义，见提取脚本），并不代表真实部署环境中的普适表现；
  \item 所提方法在本评测集上的代码提取准确率为 100\%，受益于提示架构通用模板层定义的标准化输出格式；
  \item 首次生成即通过验证的测试用例比例超过 90\%，该高首次成功率验证了分层提示架构在以最少迭代引导 LLM 生成高质量密码代码方面的有效性；
  \item \textbf{qwen2.5-coder:7b（本地、离线）}在该配置网格统计口径下 GSR/VPR 均为 100\%，FTPR 为 88.89\%；平均生成时间约 3720.70\,s，显著高于云端模型，适合对时延不敏感而对数据隐私更敏感的场景。
\end{itemize}

\subsubsection{同设置相关工作对比（DES 子集）。}
为更直接比较本文方法与代表性相关工作的表现，我们在\textbf{同一 DES 子集评测设置}（ECB/CBC/CFB/OFB × Python/C/C++，共 12 格）下整理并报告与各基线 arm 对齐的 GSR/VPR/FTPR；指标与仓库 \texttt{score} 一致（见 \texttt{相关工作对比实验.md}）。表~\ref{tab:rw_same_protocol_des} 汇总主要对比：SVEN、Self-Refine、SecCoder、AgentCoder 等基线的 FTPR 显著低于本文在 DeepSeek 上的满配结果。表中数值仍属受控桩与协议内结论，不宜外推为开放场景普适上限。

\begin{table}[H]
  \centering
  \footnotesize
  \caption{同设置相关工作对比（DES 12 格；DeepSeek 线路；GSR/VPR/FTPR 与仓库打分一致）}
  \label{tab:rw_same_protocol_des}
  \setlength{\tabcolsep}{4pt}
  \begin{tabular}{p{0.22\textwidth} p{0.38\textwidth} c c c}
    \toprule
    基线 & 设定 / arm & GSR (\%) & VPR (\%) & FTPR (\%) \\
    \midrule
    SVEN & \texttt{rw\_sven.json} / \texttt{sven\_lm\_*} & 100.0 & 0.0 & 0.0 \\
    Self-Refine+DeepSeek & \texttt{selfrefine\_deepseek} & 100.0 & 83.33 & 58.33 \\
    SecCoder+DeepSeek & \texttt{seccoder\_repro} / \texttt{seccoder\_glue\_repro} & 100.0 & 66.67 & 66.67 \\
    AgentCoder & \texttt{agentcoder\_des} & 100.0 & 41.67 & 0.0 \\
    本文方法+DeepSeek & --- & 100.0 & 100.0 & 100.0 \\
    \bottomrule
  \end{tabular}
\end{table}



\subsubsection{代码生成与提示优化中的挑战。}
\label{subsec:challenges_prompt}
一项主要发现是：\textbf{达到并维持高 FTPR 需要直面反复出现的生成失败并迭代精炼提示}——这是工作中最耗时且技术难度最高的部分。挑战出现在两个阶段：\emph{初始生成}与\emph{测试失败驱动的改进}。第~\ref{subsec:challenges_prompt_table}~节已从提示角度归纳各挑战与对策；此处从实验过程角度概括要点。

\textbf{反复出现的失败模式。}在算法-语言-LLM 组合中我们反复观察到：（i）C/C++ 使用 \texttt{uint32\_t}/\texttt{uint8\_t} 却未 \texttt{\#include <stdint.h>}，导致编译失败；（ii）\texttt{main()} 调用 \texttt{hex\_to\_bytes} 或 \texttt{sm4\_keyexp} 但同文件无定义，链接报 ``undefined reference''；（iii）生成代码将\emph{明文}十六进制当作密文输出（如 SM4-CBC）；（iv）用 \texttt{\%04x} 按 \texttt{uint16\_t} 打印密文而非按字节 \texttt{\%02x}，导致长度或字节序错误；（v）缓冲区越界与 stack smashing；（vi）因校验过严用 ``Invalid plain text'' 拒绝合法 32 位十六进制输入。此类失败仅靠测试反馈难以修复，因模型常在无显式提示约束下重复相同错误。

\textbf{提示精炼与工具。}在\textbf{生成}阶段我们向分层提示迭代加入\emph{显式}约束：强制头文件与定义\emph{顺序}（所有辅助函数在 \texttt{main()} 前）、仅按字节十六进制输出密文（\texttt{printf("\%02x", ...)}）、CBC 正确性（先与 IV 异或再块加密，禁止在输出前附加 IV）、缓冲区大小及输入校验规则（如 32 个十六进制字符合法）。当模型产生不完整或截断代码时，我们构建带 \texttt{last\_error} 与不完整重试模板的\emph{重试}提示。在\textbf{改进}阶段，我们围绕\emph{失败原因摘要}组织改进提示，优先处理编译错误并注入针对性修复指导：精确 \texttt{\#include} 映射、\texttt{goto cleanup} 修复、明文与密文及 IV 长度规则、正确 API 用法（如 \texttt{BN\_hex2bn}、\texttt{EVP\_PKEY\_encrypt}）。我们还引入轻量后处理：如在验证时若 C 代码使用 \texttt{uint32\_t} 但缺少 \texttt{stdint.h}，则在编译前自动插入。仓库中的 101 个模板体现了该迭代精炼；对特别棘手的组合（如本地 Qwen2.5-Coder 的 C 语言 SM4-CBC）我们维护专用的按 LLM、按算法提示文件。若没有这一持续的提示工程努力，报告的 FTPR 无法实现——因此我们将\textbf{失败驱动的提示优化}作为实验过程的核心贡献加以强调。


\subsubsection{错误模式分类效果。}
（口径见「消融与错误统计口径」。）对 DeepSeek 在全算法全模式网格上写入的 \texttt{llm\_performance.json} 中失败记录（校验失败或向量未通过）经启发式归入四类；共 101 条失败样本，合计「修复后通过率」（本条或同任务更晚日志中 \texttt{test\_success} 为真）为 87.13\%（表~\ref{tab:error_repair}）。实现细节类条数最多；算法理解类修复后通过率相对较低；环境配置与代码结构类在日志意义上均已跑通。主要结论如下：
\begin{enumerate}[leftmargin=1.2em, labelsep=0.4em]
  \item \textbf{分布与难点。}实现细节类失败条数最多（41，40.59\%）；算法理解类修复后通过率最低（77.14\%），仍可从模板层强化算法原理（如 OFB-64 与 OFB-8）中获益。
  \item \textbf{易收敛类型。}环境配置类与代码结构类在观测口径下修复后通过率均为 100.0\%，特征相对明确（如环境变量名、缺头文件、结构问题），针对性修复指导易见效。
  \item \textbf{整体效果。}自 0\% 基线至合计 87.13\% 的修复后通过率，验证了错误模式分类与日志聚合口径下针对性修复的有效性（详见表注与仓库聚合脚本说明）。
\end{enumerate}

\begin{table}[H]
  \centering
  \footnotesize
  \caption{不同错误类型的分布与修复效果（DeepSeek；四算法全支持模式$\times$三语言；\texttt{llm\_performance.json} 聚合）}
  \label{tab:error_repair}
  \setlength{\tabcolsep}{3pt}
  \begin{tabular}{
    p{0.28\textwidth}                       
    >{\centering\arraybackslash}p{0.08\textwidth} 
    >{\centering\arraybackslash}p{0.12\textwidth} 
    >{\centering\arraybackslash}p{0.16\textwidth}  
    >{\centering\arraybackslash}p{0.18\textwidth}  
  }
    \toprule
    错误类型               & 数量 & 占比 & 修复前通过率 & 修复后通过率 \\
    \midrule
    算法理解  & 35    & 34.65\%    & 0.0\%                & 77.14\%               \\
    实现细节  & 41    & 40.59\%    & 0.0\%                & 87.80\%               \\
    环境配置  & 12    & 11.88\%    & 0.0\%                & 100.0\%               \\
    代码结构  & 13    & 12.87\%    & 0.0\%                & 100.0\%               \\
    \midrule
    合计      & 101   & 100.0\%    & 0.0\%                & 87.13\%               \\
    \bottomrule
  \end{tabular}
\end{table}

\subsection{消融实验}
通过消融实验量化所提方法各核心模块的贡献。表~\ref{tab:ablation_main} 给出主要消融实验结果，关键结论如下。

\begin{table}[H]
  \centering
  \footnotesize
  \caption{主要消融实验结果（DeepSeek、豆包与 Claude；四算法全支持模式$\times$三语言）}
  \label{tab:ablation_main}
  \setlength{\tabcolsep}{4pt}
  \begin{tabular}{l c c c}
    \toprule
    系统配置                  & VPR (\%) & FTPR (\%) & 性能下降 (VPR/FTPR) \\
    \midrule
    完整所提方法                  & 100.0    & 100.0     & -- / --                     \\
    无测试反馈改进         & 100.0    & 77.08     & 0.0\% / -22.92\%            \\
    无分层提示架构（仅 base_prompt.yaml）  & 60.42    & 50.0      & -39.58\% / -50.0\%          \\
    无任何提示  & 0.0    & 0.0     & -100.0\% / -100.0\%         \\
    \bottomrule
  \end{tabular}
\end{table}

\begin{enumerate}[leftmargin=1.2em, labelsep=0.4em]
  \item \textbf{测试反馈改进主要作用于 FTPR}：去除闭环后 FTPR 降至 77.08\%（表~\ref{tab:ablation_main}），说明\textbf{向量一致性功能正确性仍依赖测试反馈}。
  \item \textbf{分层提示架构同时抬高 VPR 与 FTPR}：仅 \texttt{common\_only} 时 VPR/FTPR 为 60.42\%/50.0\%。
  \item \textbf{完全撤除领域栈则指标塌缩}：\texttt{no\_prompt} 下仅保留单行任务描述，表中 VPR/FTPR 仍为 0\%。
\end{enumerate}


\subsubsection{提示优化消融。}
针对提示策略的进一步消融验证了分层可组合设计的价值。表~\ref{tab:ablation_prompt} 给出不同提示策略的对比结果：
\begin{itemize}[leftmargin=1.2em, labelsep=0.4em]
  \item 随\textbf{提示模板层次}由通用向「通用+LLM+语言」递进，VPR/FTPR 逐步升高（表~\ref{tab:ablation_prompt}），说明组合多层模板能充分发挥领域知识与模型特性；
  \item \textbf{四层可组合提示策略（所提方法）}达到最佳性能（FTPR 100.0\%），显著优于单一通用模板（25.0\%）及中间档位（79.17\%--91.67\%），证实分层提示工程能有效弥补通用 LLM 在密码学知识上的不足。
\end{itemize}

\begin{table}[H]
  \centering
  %\footnotesize
  \caption{提示策略消融实验结果（DeepSeek、豆包与 Claude；四算法全支持模式$\times$三语言）}
  \label{tab:ablation_prompt}
  \setlength{\tabcolsep}{4pt}
  \begin{tabular}{l c c c}
    \toprule
    提示策略               & GSR (\%) & VPR (\%) & FTPR (\%) \\
    \midrule
    单一通用模板         & 100.0    & 45.83    & 25.0      \\
    通用+LLM              & 100.0    & 87.5     & 79.17     \\
    通用+LLM+语言        & 100.0    & 91.67    & 91.67     \\
    通用 +LLM+ 语言+算法（本文方法） & 100.0 & 100.0 & 100.0 \\
    \bottomrule
  \end{tabular}
\end{table}

\section{结论与未来工作}
\subsection{研究结论}
针对传统 LLM 在密码学代码生成中错误率高、缺乏领域适配的问题，本文提出一种基于测试反馈与分层提示的方法：建立闭环测试反馈系统以实现自动错误识别与针对性修复；设计分层可组合的领域提示工程，并采用\textbf{失败驱动的提示精炼}（迭代应对反复出现的生成失败）；建立密码学代码错误模式分类系统。我们强调，持续、失败驱动的提示优化是实验过程中要求最高、最耗时的部分，对报告的功能通过率至关重要。

\textbf{实验验证}：在 159 个真实测试用例中，\emph{代码生成与提取成功率}（生成可运行代码且正确解析）在本基准上达到 100\%；\emph{功能测试通过率}（输出与 NIST/FIPS/GM 标准向量一致）在完整方法下由初始 45.1\% 提升至 93.33\%，并明显优于通用单提示基线（与表~\ref{tab:ablation_prompt} 相同网格：DeepSeek、豆包与 Claude ×全算法全模式，仅 25.0\% FTPR）。其中，100\% 来自历史聚合口径；93.33\% 来自固定批次主实验口径，两者不可直接横比。上述结论均限定于本文基准、固定测试桩与有限迭代预算，不构成对开放真实工程场景的普适正确性保证。

本文贡献体现在理论、技术与实践三个层面。
\begin{itemize}
    \item \textbf{理论层面}：提出由错误模式分类驱动的测试反馈闭环机制，构建领域分层提示工程范式，设计离线大模型本地化引擎，为密码学代码生成提供理论框架与实现方案。
    \item \textbf{技术层面}：实现高精度代码提取；失败驱动的提示精炼应对反复出现的生成失败（如强制头文件、函数顺序、输出格式）；错误匹配与修复策略在 \texttt{llm\_performance.json} 失败样本聚合口径下使整体修复后通过率达 87.13\%。
    \item \textbf{实践层面}：建立 101 个提示模板库与 159 个真实用例数据集；开发具备多 LLM 切换、批量生成等功能的完整系统，可直接落地使用。
\end{itemize}

\subsection{研究局限}
本文方法仍存在以下局限：
\begin{itemize}
    \item \textbf{错误模式覆盖不完整}。当前错误模式分类系统基于 159 个测试用例构建，可能尚未覆盖部分类型的密码学代码错误，完整性有待提升。表~\ref{tab:error_repair} 的量化统计仅反映 DeepSeek 在全网格上的日志口径，未必推广到其他提供商。在当前该聚合中，算法理解类修复后通过率相对较低（77.14\%），复杂算法逻辑错误的纠正能力仍需优化。
    \item \textbf{算法支持有限}。当前方法主要支持 DES、AES、SM4、RSA 四种算法，尚不支持椭圆曲线加密、ChaCha20 等，也未覆盖后量子密码、同态加密等新兴算法。
    \item \textbf{安全检测维度不足}。当前方法主要关注代码功能正确性，尚未深入侧信道攻击、时序攻击等安全漏洞检测。
    \item \textbf{构念与外部有效性边界}。FTPR 衡量的是“标准向量上的功能正确性”，不等同于性能、常数时间性质、侧信道鲁棒性与工程可维护性；在 159 例受控基准上的优势不自动外推到复杂生产环境。
\end{itemize}

\subsection{未来工作方向}
针对上述局限，未来将开展以下工作：
\begin{itemize}
    \item \textbf{扩展错误模式库}。收集更多密码学代码生成的测试用例，扩展错误模式分类体系，提升错误识别准确率；
    \item \textbf{支持新兴密码算法}。为新兴算法增加更多提示模板，扩大方法适用范围；
    \item \textbf{增强安全检测能力}。结合形式化验证、静态分析等技术，检测代码中的安全漏洞，提升密码代码安全性；
    \item \textbf{实现在线学习能力}，根据用户反馈持续优化错误模式库与修复策略，增强系统自适应能力。
\end{itemize}

% \section*{Acknowledgements}
% This work was supported by Sichuan University. The authors thank the teachers and classmates of the School of Cyber Science and Engineering of Sichuan University for their valuable suggestions and discussions. Doubao was used to produce parts of the work, including text, tables, charts, codes, data, and quotations.

\renewcommand{\refname}{参考文献}
\renewcommand{\bibname}{参考文献}
\bibliographystyle{splncs04}
\bibliography{refs} 
%% Appendix (if needed)
%\appendix
%\section{Appendix Title}
%Appendix content here.

\end{document}